\section{Mapa de lectura: de EP137 a los episodios siguientes de la serie}

EP137 ocupa un lugar particular en la serie de IA de Zhang Xiaojun. EP140 y EP138 se centran más en los laboratorios de modelos y en el paradigma industrial de los Agent; EP139 ofrece una historia técnica de los Agent; EP137, en cambio, lleva la cuestión hacia la investigación científica y la demostración formal. Es el puente entre «la IA transforma la productividad» y «la IA transforma la producción de conocimiento».

\subsection{Cinco preguntas que el lector debería llevarse}

Primera: ¿puede la IA pasar de resolver problemas a hacer verdadera investigación matemática? Segunda: ¿se convertirá Lean en la interfaz principal de AI for Math? Tercera: ¿puede un modelo aprender la intuición de los matemáticos, o esta tendrá que aportarla el ser humano a largo plazo? Cuarta: ¿cómo puede un Neo Lab como Axiom encontrar productos intermedios entre la investigación básica y la comercialización? Quinta: ¿el éxito de AI for Math impulsará, a su vez, el avance de los algoritmos y de la teoría de la propia IA?

\begin{knowledgebox}{Relación entre este episodio y las notas posteriores}
Si más adelante se elaboran notas de otras entrevistas sobre AI for Science, robótica, modelos del mundo o conducción autónoma, puede reutilizarse el marco de este episodio: primero buscar un entorno verificable, después examinar la retroalimentación de los expertos humanos y, por último, ver si el modelo puede formar un ciclo cerrado de largo plazo. Las matemáticas son solo el tipo de escenario de este marco en el que la verificación es más rigurosa.
\end{knowledgebox}

\subsection{Resumen de la sección}

El valor de EP137 está en que amplía los límites de la serie de IA: de los productos Agent y el entrenamiento de modelos pasa a las matemáticas y al descubrimiento científico. Nos recuerda que las aplicaciones de IA más difíciles suelen requerir, al mismo tiempo, verificación formal y sentido estético humano.
